% Exact table extracted from the pinned manuscript; needs booktabs/amsmath.
\begin{table}[!htbp]
\centering
\caption{Selected model configurations and common evaluation protocol.}
\label{tab:Table03_Training_And_RFE_Protocol}
\small
\setlength{\tabcolsep}{4pt}
\renewcommand{\arraystretch}{1.15}
\begin{tabular}{@{}l l@{}}
\toprule
Model & Key configuration \\
\midrule
TabM & 3 blocks, 32 members, dropout 0.1, AdamW \\
MLP & 2 hidden layers (64, 32), dropout 0.2 \\
BiLSTM & Hidden 32 per direction, head dropout 0.2 \\
GCN & 2 layers, gated fusion, head dropout 0.3 \\
Transformer & 2 encoder layers, 4 heads, dropout 0.1 \\
Bai LSTM & LSTM 32, dense 16--32--64--1, RMSprop \\
SVR-RBF & $C=1$, $\epsilon=0.1$, RBF $\gamma=\mathrm{scale}$ \\
CatBoost & Depth 6, lr 0.05, leaf L2 reg 10 \\
Random Forest & Depth 16, min leaf 5, selected 7-tree prefix \\
Elastic Net & $\alpha=0.01$, L1 ratio 0.75 \\
\bottomrule
\end{tabular}
\par\smallskip
\parbox{\columnwidth}{\footnotesize
\textbf{Note:}
During training, MLP, GCN, CatBoost, and Random Forest
assign greater weights to rainy and non-rainy overcast
hours. The remaining models give each training hour
equal weight.
TabM trains for 54 epochs
(lr $2\times10^{-3}$, weight decay 0.01).
SVR-RBF and Elastic Net run until solver convergence.
}
\end{table}
